\chapter{Communicatie via hormonen}\label{ch:g8:hormonal-communication}

De \omterm{def:g8:puberty:puberty}{puberteit} stelde ze voor
(\cref{def:g8:puberty:hormones}); de pil buit ze uit
(\cref{def:g8:contraception:pill}); en de hete golf van schrik bezorgde er
één midden in een sprint
(\cref{ex:g8:nervous-communication:cases}). De postdienst verdient haar
eigen rondgang: de \omterm{def:g8:hormonal-communication:glands}{klieren} die de brieven schrijven, het \omterm{prop:g4:journey-of-food:absorb}{bloed} dat ze
draagt, de adresseringstruc die van een omroep iets precies maakt --- en
de \omterm{prop:g8:hormonal-communication:feedback}{terugkoppelingen} die de hele briefwisseling in evenwicht houden.

\section{Klieren, hormonen, doelwitten}

\begin{definition}[De hormoonklieren]\label{def:g8:hormonal-communication:glands}
Een \emph{klier}\index{klier} waarvan het product \emph{in het \omterm{prop:g4:journey-of-food:absorb}{bloed}}
wordt afgegeven (en niet door een buisje, zoals bij het \omterm{def:g4:journey-of-food:tract}{speeksel}) is een
\emph{hormoonklier}, en haar producten zijn de
\emph{hormonen}\index{hormoon}. De belangrijkste schrijvers: de
\emph{regelklier} onder in de \omterm{def:g5:brain-in-command:system}{hersenen} --- de starter van de \omterm{def:g8:puberty:puberty}{puberteit}, en
de dirigent van de meeste andere; de \emph{schildklier} in de \omterm{def:g1:my-body:parts}{hals} ---
haar hormoon bepaalt het draaitempo van het hele lichaam; de
\emph{bijnieren} als kapjes op de \omterm{def:g7:removing-wastes:kidneys}{nieren} --- het alarmhormoon van schrik en
sprint; de \emph{alvleesklier} --- haar hormoon beheert de \omterm{prop:g7:organs-at-work:bill}{glucose} in het
\omterm{prop:g4:journey-of-food:absorb}{bloed}; en de \emph{\omterm{def:g8:reproductive-systems:female}{eierstokken} en \omterm{def:g8:reproductive-systems:male}{zaadballen}} --- de cyclus- en
verbouwingshormonen uit \cref{ch:g8:reproductive-systems}.
\end{definition}

\begin{proposition}[Adresseren via receptoren]\label{prop:g8:hormonal-communication:address}
Het \omterm{prop:g4:journey-of-food:absorb}{bloed} bezorgt elk \omterm{def:g8:hormonal-communication:glands}{hormoon} overal
(\cref{ex:g8:puberty:sites}) --- en toch werkt elk \omterm{def:g8:hormonal-communication:glands}{hormoon} alleen in op
zijn \emph{doelorganen}\index{doelorgaan}: de cellen die met
aanmeerplaatsen, \emph{receptoren}, voor die boodschapper gebouwd zijn. Het
vervoer is een omroep, de adressering zit bij de ontvanger: een brief die
elk huis ontvangt, maar die alleen leesbaar is waar de sleutel past. (Het
aanmeren bij de \omterm{def:g8:nervous-communication:synapse}{synaps} uit
\cref{def:g8:nervous-communication:synapse} is dezelfde truc op
nanometerafstand.)
\end{proposition}
\admitted

\begin{example}[De rondgang van één boodschapper: het alarmhormoon]\label{ex:g8:hormonal-communication:alarm}
Een bijna-aanrijding op de fiets: de bijnieren storten hun alarmhormoon in
het \omterm{prop:g4:journey-of-food:absorb}{bloed}, en binnen enkele seconden --- één ronde van de bloedsomloop
(\cref{ex:g4:heart-and-blood:round}) --- reageert elke passende ontvanger
tegelijk: het tempo van het \omterm{def:g4:heart-and-blood:heart}{hart} springt omhoog, de luchtwegen worden
wijder, de \omterm{prop:g7:digestion-and-nutrients:absorption}{lever} geeft gebankierde \omterm{prop:g7:organs-at-work:bill}{glucose} vrij (de bank uit
\cref{ex:g7:digestion-and-nutrients:breadtrace}, opgeëist), de vaten van de
\omterm{def:g3:bones-and-muscles:muscle}{spieren} worden wijder en die van het \omterm{def:g4:journey-of-food:tract}{spijsverteringskanaal} nauwer (het
omleiden uit \cref{prop:g7:organs-at-work:supply}, nu door een \omterm{def:g8:hormonal-communication:glands}{hormoon}
vastgehouden), en de \omterm{def:g1:the-five-senses:sense}{huid} wordt bleek en gaat prikkelen. Eén brief, zes
antwoorden, één gecoördineerde toestand: klaar.
\end{example}

\section{Evenwicht door terugkoppeling}

\begin{proposition}[Het terugkoppelingsprincipe]\label{prop:g8:hormonal-communication:feedback}
Hormoonstelsels regelen zichzelf met
\emph{terugkoppeling}\index{terugkoppeling}: de schrijvende \omterm{def:g8:hormonal-communication:glands}{klier}
\emph{houdt het resultaat van haar eigen brieven in de gaten} en
vertraagt zodra dat resultaat bereikt is --- een thermostaat, geen kraan.
Het paradepaardje is de \omterm{prop:g7:organs-at-work:bill}{glucose} in het \omterm{prop:g4:journey-of-food:absorb}{bloed}: na een maaltijd zet de
stijgende \omterm{prop:g7:organs-at-work:bill}{glucose} zelf het opslaghormoon van de alvleesklier aan; als de
\omterm{prop:g7:organs-at-work:bill}{glucose} gebankierd is, zakt het gehalte, en de afgifte zakt mee; en tussen
de maaltijden door maakt een partnerhormoon de bank de andere kant op weer
los. Het gehalte rijdt dag en nacht in een smalle baan, zonder bewuste
bestuurder.
\end{proposition}
\admitted

\begin{example}[Als de terugkoppeling faalt: diabetes]\label{ex:g8:hormonal-communication:diabetes}
Bij \emph{diabetes}\index{diabetes} valt het signaal van het opslaghormoon
uit --- de alvleesklier stopt met schrijven, of de ontvangers stoppen met
antwoorden. Na elke maaltijd stapelt de \omterm{prop:g7:organs-at-work:bill}{glucose} zich in het \omterm{prop:g4:journey-of-food:absorb}{bloed} op; het
terugwinnen door de \omterm{def:g7:removing-wastes:kidneys}{nier} wordt overweldigd en er verschijnt suiker in de
\omterm{def:g7:removing-wastes:kidneys}{urine} --- het klassieke signaal uit
\cref{ex:g7:removing-wastes:thrift}, eindelijk verklaard. De behandeling
is toegepast hoofdstuk: afgemeten doses van het ontbrekende \omterm{def:g8:hormonal-communication:glands}{hormoon}, getimed
op de maaltijden --- de brief per injectiespuit geschreven. Miljoenen mensen
leven prima met precies die regeling.
\end{example}

\begin{remark}[De twee stelsels in partnerschap]\label{rem:g8:hormonal-communication:partnership}
De telegraaf en de post uit
\cref{prop:g8:nervous-communication:compare} vormen één bestuur. De
regelklier onder in de \omterm{def:g5:brain-in-command:system}{hersenen} is de koppeling: het \omterm{def:g5:brain-in-command:system}{zenuwstelsel} erboven,
de \omterm{def:g8:hormonal-communication:glands}{hormonen} eronder --- de schrik van de fietser liep eerst via de \omterm{def:g5:brain-in-command:system}{zenuwen}
en werd door \omterm{def:g8:hormonal-communication:glands}{hormonen} volgehouden; de \omterm{def:g8:puberty:puberty}{puberteit} loopt traag op \omterm{def:g8:hormonal-communication:glands}{hormonen}
onder een dirigent in de \omterm{def:g5:brain-in-command:system}{hersenen}; en de schapen uit
\cref{ex:g8:reproduction-and-environment:lambs} zetten daglengte (een
waarneming via de \omterm{def:g5:brain-in-command:system}{zenuwen}) om in voortplantingshormonen. Snelle draden waar
snelheid telt, trage post waar het programma maanden moet lopen, en
vertaalkantoren daartussen.
\end{remark}

\section{Een hormoonstelsel lezen}

\begin{method}[De vier vragen bij elk hormoon]\label{met:g8:hormonal-communication:read}
Stel bij elk \omterm{def:g8:hormonal-communication:glands}{hormoon} dat je tegenkomt vast:
\begin{enumerate}
  \item \emph{schrijver}: welke \omterm{def:g8:hormonal-communication:glands}{klier}, in het \omterm{prop:g4:journey-of-food:absorb}{bloed}, op welk sein?
  \item \emph{ontvangers}: welke \omterm{prop:g8:hormonal-communication:address}{doelorganen} dragen zijn receptoren, en
        wat doet elk bij bezorging?
  \item \emph{\omterm{prop:g8:hormonal-communication:feedback}{terugkoppeling}}: welk resultaat meldt zich terug, en hoe
        wordt te veel schrijven voorkomen?
  \item \emph{falen en gebruik}: wat gebeurt er als de brief uitblijft ---
        en bootst de geneeskunde hem na (de pil, de doses van de
        diabetespatiënt)?
\end{enumerate}
\end{method}

\begin{example}[De methode op de tempoklier]\label{ex:g8:hormonal-communication:thyroid}
Schrijver: de schildklier, gedirigeerd door de regelklier. Ontvangers:
vrijwel elke \omterm{def:g5:first-look-at-cells:cell}{cel} --- het \omterm{def:g8:hormonal-communication:glands}{hormoon} bepaalt het tempo van het trage vuur uit
\cref{prop:g7:what-is-respiration:power}. \omterm{prop:g8:hormonal-communication:feedback}{Terugkoppeling}: de regelklier
leest het gehalte af en stelt haar dirigeren bij. Falen: te weinig, en de
hele motor draait stationair --- moeheid, kou, alles vertraagd; te veel, en
hij raast. Gebruik: het ontbrekende \omterm{def:g8:hormonal-communication:glands}{hormoon} als dagelijkse tablet --- een
van de oudste brieven-per-post van de geneeskunde.
\end{example}

\section{Oefeningen}

\begin{exercise}[$\star$]\label{exo:g8:hormonal-communication:1}
Wat maakt een \omterm{def:g8:hormonal-communication:glands}{klier} een hormoonklier? Noem er vier, met van elk één rol van
haar \omterm{def:g8:hormonal-communication:glands}{hormoon}.
\end{exercise}

\begin{exercise}[$\star$]\label{exo:g8:hormonal-communication:2}
Het \omterm{prop:g4:journey-of-food:absorb}{bloed} bezorgt overal; \omterm{def:g8:hormonal-communication:glands}{hormonen} werken ergens. Los die schijnbare
tegenstelling op.
\end{exercise}

\begin{exercise}[$\star$]\label{exo:g8:hormonal-communication:3}
Noem vier van de zes antwoorden op het alarmhormoon, elk met het
hoofdstuk dat erover ging.
\end{exercise}

\begin{exercise}[$\star$]\label{exo:g8:hormonal-communication:4}
Beschrijf het terugkoppelingsprincipe, en zijn paradepaardje.
\end{exercise}

\begin{exercise}[$\star$]\label{exo:g8:hormonal-communication:5}
Wat faalt er bij \omterm{ex:g8:hormonal-communication:diabetes}{diabetes}, en wat verschijnt er daardoor opnieuw bij de
\omterm{def:g7:removing-wastes:kidneys}{nier}?
\end{exercise}

\begin{exercise}[$\star$]\label{exo:g8:hormonal-communication:6}
Waar raken het \omterm{def:g5:brain-in-command:system}{zenuwstelsel} en het hormoonstelsel elkaar? Geef één
voorbeeld van een estafette.
\end{exercise}

\begin{exercise}[$\star\star$]\label{exo:g8:hormonal-communication:7}
Waarom is een thermostaat het juiste beeld voor \omterm{prop:g8:hormonal-communication:feedback}{terugkoppeling}, en een
kraan het verkeerde?
\end{exercise}

\begin{exercise}[$\star\star$]\label{exo:g8:hormonal-communication:8}
Laat \cref{met:g8:hormonal-communication:read} lopen over het
opslaghormoon van de alvleesklier.
\end{exercise}

\begin{exercise}[$\star\star$]\label{exo:g8:hormonal-communication:9}
Leg uit dat de injectiespuit van de diabetespatiënt en de \omterm{def:g8:contraception:pill}{anticonceptiepil}
dezelfde \omterm{def:g5:biodiversity:species}{soort} handeling zijn, met vraag 4 van de methode.
\end{exercise}

\begin{exercise}[$\star\star$]\label{exo:g8:hormonal-communication:10}
Wijs de stelsels toe, met redenen: het begin van het rillen; het groeien
van een baard; de smalle baan van de bloedglucose; de kniepeesreflex.
\end{exercise}

\begin{exercise}[$\star\star$]\label{exo:g8:hormonal-communication:11}
De schapen uit \cref{ex:g8:reproduction-and-environment:lambs} zetten licht
om in lammeren. Volg de estafette door de kantoren van
\cref{rem:g8:hormonal-communication:partnership}.
\end{exercise}

\begin{exercise}[$\star\star\star$]\label{exo:g8:hormonal-communication:12}
``Receptoren maken van een omroep een net van privélijnen.'' Verdedig die
zin, en zet haar dan door: wat moet er in een \omterm{def:g5:first-look-at-cells:cell}{cel} veranderen om zich bij
het publiek van een \omterm{def:g8:hormonal-communication:glands}{hormoon} te \emph{voegen} --- en wat suggereert dat over
hoe de doelwitten van de \omterm{def:g8:puberty:puberty}{puberteit} ooit zijn gaan luisteren?
\end{exercise}

\section{Probleem: Het glucoselogboek}

\begin{problem}[{Weekendprobleem --- een dag bloedsuiker, gelezen als
briefwisseling}]\label{pb:g8:hormonal-communication:1}
Het (verzonnen) dagverslag van een continue glucosemeter: 's nachts
stabiel; ontbijt om 8 uur --- een stijging, daarna een rustige daling tot
10 uur; sport om 12 uur --- een kort dipje, snel opgevangen; een
overgeslagen lunch --- de hele middag toch stabiel; een schrik om 17 uur
(een bijna-aanrijding op de fiets) --- een scherpe korte stijging; avondeten
om 19 uur --- opnieuw stijging en daling.

\medskip\noindent\textbf{Deel I --- De ochtend.}
\begin{enumerate}
  \item De stabiliteit 's nachts: welke briefwisseling van twee brieven
        hield de baan vast terwijl haar eigenaar sliep?
  \item De stijging en daling na het ontbijt: noem het sein, de schrijver,
        de brief, en de ontvangers die het overschot bankierden.
  \item Waarom is het \emph{stoppen} van de daling --- geen duik onder de
        baan --- op zichzelf de handtekening van de \omterm{prop:g8:hormonal-communication:feedback}{terugkoppeling}?
  \item Het gehalte om 10 uur is gelijk aan dat om 7 uur. Wat voorspelt
        het beeld van de thermostaat over het verloop van de afgifte
        daartussen?
\end{enumerate}

\medskip\noindent\textbf{Deel II --- De middag.}
\begin{enumerate}[resume]
  \item Het dipje bij het sporten en het opvangen ervan: welke brief riep
        de bank van de \omterm{prop:g7:digestion-and-nutrients:absorption}{lever} aan, en de rekening uit welk hoofdstuk werd
        er uitgegeven?
  \item De overgeslagen lunch veranderde niets zichtbaars. Het stille werk
        van welk partnerhormoon overbrugde de middag?
  \item De scherpe stijging bij de schrik: noem de schrijver en het punt
        van die plotselinge suiker --- klaar waarvoor, volgens
        \cref{ex:g8:hormonal-communication:alarm}?
  \item Waarom zakte de stijging van de schrik binnen het uur weer weg,
        anders dan de blijvende opeenhoping bij \omterm{ex:g8:hormonal-communication:diabetes}{diabetes}?
\end{enumerate}

\medskip\noindent\textbf{Deel III --- Het consult.}
\begin{enumerate}[resume]
  \item Het logboek van een diabetespatiënt zou van diezelfde dag bij elke
        maaltijd afwijken. Beschrijf de onbehandelde curve na het ontbijt,
        en waar het signaal van de \omterm{def:g7:removing-wastes:kidneys}{nier} zou opduiken.
  \item De behandelde versie: wat bootst de afgemeten dosis vóór elke
        maaltijd na, en waarom moet haar grootte bij het bord passen?
  \item De arts noemt het gezonde logboek ``een gesprek dat je nooit
        hoort''. Noem de sprekers van die dag --- \omterm{def:g8:hormonal-communication:glands}{klieren}, brieven,
        ontvangers --- in één samenvattende tabel of alinea.
  \item Sluit af met de moraal van het hoofdstuk in één zin: wat voor
        stelsel houdt een waarde een leven lang in een baan zonder
        bestuurder?
\end{enumerate}
\end{problem}
